% Chinese production target: zh-Hans-CN mathematical register.
% Source authority: sealed German P31 cumulative TeX, SHA-256
% A48CB5CD1716974B686AC1CBA681CA4B17BC72F9043B78AD2528ACA41FCF814F.
% Exact logical German article SHA-256:
% A61962BD0D031352C9EB06FDD7A39813B49289B8361E22D5082F49B0AB0BD439.
% This file is not a Singapore, Taiwan, Hong Kong, or Macao localization.
\documentclass[11pt]{article}
\usepackage[a4paper,margin=2.35cm]{geometry}
\usepackage{fontspec}
\usepackage{xeCJK}
\usepackage{amsmath,amssymb}
\setmainfont{Noto Serif}
\setCJKmainfont[BoldFont=Microsoft YaHei]{SimSun}
\setCJKsansfont{Microsoft YaHei}
\XeTeXlinebreaklocale "zh"
\XeTeXlinebreakskip=0pt plus 1pt
\setlength{\parindent}{2em}
\setlength{\parskip}{0.45em}
\setlength{\emergencystretch}{3em}
\newcommand{\srcemph}[1]{\textbf{#1}}
\newcommand{\srctextsc}[1]{\textbf{#1}}
\begin{document}

\setcounter{footnote}{0}
\section*{33. 算術觀點下的超復量與表示論}
\srcemph{Atti Congresso Bologna 2 (1928), S. 71--73.}

\begingroup
\renewcommand{\thefootnote}{(\arabic{footnote})}

\begin{center}
\textsc{Emmy Noether} (Göttingen -- Germania)\\[-0.1em]
\rule{3em}{0.4pt}\\[0.8em]
\srctextsc{超復量與表示論}\\
\srctextsc{（算術觀點下）}\footnote{詳細論述將發表於 Math. Zeitschr.；在那裏的一篇續文還將討論非交換伽羅瓦理論。關於“分裂域”，亦參見 R. Brauer 與 E. Noether, Ber. Berl. Ak. 1927；關於一般結構定理，參見 G. Köthe 在本卷 Atti 中的報告。}
\end{center}

我想說明，怎樣把表示論——尤其是羣表示論——理解爲模的算子同構類與理想的算子同構類的理論，以及怎樣把後兩者納入一個推廣的羣概念，即帶算子的羣。\footnote{帶算子的羣可追溯到 W. Krull (Math. Zeitschr. 23) 和 O. Schmidt (Math. Zeitschr. 29)。本文對給定表示之約化問題的處理，在表示除環爲交換的限制下，也來自 W. Krull, Theorie und Anwendung der verallgemeinerten Abelschen Gruppen, Heidelberger Ber. 1926。}

所謂\srcemph{一個羣的表示}——在預先給定的除環 $P$ 中——衆所周知，是指如下情形：對元素 $a,b,\ldots$（它們屬於羣 $\mathfrak{G}$）分別配以矩陣 $A,B,\ldots$，矩陣的係數取自 $P$，使乘積 $ab$ 對應於乘積 $AB$；於是這個矩陣系統成爲該羣的一個同態像。與此同時，也給出了“羣環”的一個表示；羣環由所有“羣數”（即羣環元素）組成，也就是所有線性組合 $a\alpha+b\beta+\cdots+c\gamma$，其中 $a,b,\ldots,c$ 每次取有限多個羣元素，而 $\alpha,\beta,\ldots,\gamma$ 相應地取 $P$ 中有限多個元素；這裏把 $a,b,\ldots,c$ 視爲線性無關，乘法由羣乘法和運算法則定義。羣環的表示同時關於加法同態；因此，羣的表示被納入\srcemph{環的表示}這一一般問題之下，在其中要求同時關於加法和乘法同態。

這裏有兩個本質問題。第一個是\srcemph{給定表示的約化}問題，以及在約化中出現的不可約組成部分是否唯一的問題；第二個則是詢問一個給定環在一個給定的、不必交換的除環中的全部可能表示，例如只考慮其中的不可約表示。

第一個問題遠爲簡單；這一點已經表現在，對於待表示的環以及通常非交換的表示除環，都不需要任何特殊假設。表示由固定階數的矩陣給出這一事實，已經提供了一切必要的有限性假設。藉助帶算子的羣的理論，這個問題可以完全處理；因此我先簡短說明這一理論。

稱羣 $\mathfrak{B}$ 爲一個\srcemph{帶算子系的羣}，如果同時給定一個符號系統 $\Theta,H,\ldots$——即算子系——並且表達式 $\Theta(a),H(a),\ldots$ 對每個 $a$（取自 $\mathfrak{B}$）都唯一確定，且給出 $\mathfrak{B}$ 中的一個元素，而且滿足分配律 $\Theta(ab)=\Theta(a)\Theta(b)$。兩個帶有同一算子系的羣稱爲算子同構，如果它們在通常意義下同構，並且由 $a$ 與 $\bar a$ 的對應還推出 $\Theta(a)$ 與 $\Theta(\bar a)$、$H(a)$ 與 $H(\bar a),\ldots$ 的對應。若只把自身也容許該算子系的子羣算作可容許子羣（因而所謂單羣，是指除單位元外沒有可容許的真正規子羣），那麼關於合成列的若爾當--赫爾德定理仍以如下形式成立：如果一個帶算子系的羣確有合成列，那麼其因子羣（合成因子）的系統在算子同構意義下唯一確定，至多差一個排列。在同一假設下，分解爲直接不可分解因子的定理也成立，並且這種分解在算子同構意義下唯一。

它同表示問題的聯繫在於，理想和模尤其屬於帶算子的羣：它們相對於加法被看作阿貝爾羣，而算子由同環元素相乘給出。每一個表示都由一個表示模產生，也就是由一個關於該環和表示除環的模產生。更準確地說，每一個模類，即彼此算子同構的表示模所成的類，產生同一個表示。這樣，表示約化的唯一性問題便由合成列定理和直積定理回答。把合成列應用於表示模，就對每個矩陣給出如下形式的約化：
\[
\begin{pmatrix}
B_1&0&0\\
&B_2&0\\
A_{ik}&\cdots&B_r
\end{pmatrix},
\]
其中對應於因子羣的對角矩陣產生一個“不可約”表示，也就是說，它本身不再容許這樣的約化。

既然這些因子羣的類唯一確定，對角表示也同樣唯一，這裏的唯一性是在表示等價的意義下。類似地，直積定理給出“分解”爲“不可分解”組成部分的唯一性：
\[
\begin{pmatrix}
C_1&&&\\
&C_2&&\\
&&\cdots&\\
&&&C_s
\end{pmatrix},
\]
其中每個 $C_i$ 隨後還可以按合成列定理唯一地繼續約化。

我想在超復系統的表示情形下，更一般地在滿足“雙鏈條件”的環的情形下，勾畫\srcemph{第二個問題}。其結果是，每一個不可約（單）模的算子同構類都是某個理想的算子同構類；因而，由單側理想組成的合成列已經通過其合成因子給出全部不可約表示。更準確地說，只須考察模去根基（在此即 Jacobson 根）所得商環中的合成列，而該商環成爲一個“完全可約”環（現代稱半單環）。問題由此化歸爲完全可約環的結構研究；因而，把自同態除環中的表示作爲出發點就顯得與問題相適應。這裏所說的是如下事實：這樣一個環的一個雙側單分量中的所有單右理想都屬於同一個類，因此有同一個自同態環；由於這些理想是單的，這個自同態環成爲除環，同時也是單左理想的自同態環。這個單環本身同自同態除環上的全體矩陣系統同構；而從這一表示出發，便產生關於一個子除環的一切表示，特別是關於該超復系統的係數域的一切表示，只要把自同態除環的元素本身替換爲相應的矩陣。於是，在超復系統情形下已知的 Wedderburn 結構定理，在這裏通過來自一般羣論的自同態除環概念得到新的解釋，並且同時被識別爲超復系統表示論的來源。由此產生非交換除環的伽羅瓦理論中的新問題，它們同“分裂域”問題緊密相連。同時，通過不可分解理想的自同態環，也展望到一般的、並非完全可約的環的結構定理。羣論，在它擴展到最一般的帶算子的羣之後，在這裏也再次顯現爲一種整理原則。

\endgroup

\end{document}
